\documentclass[10pt,a4paper]{amsart}
\usepackage[margin=19mm]{geometry}
\usepackage{amsmath,amssymb,mathtools}
\usepackage[scr=rsfs,cal=euler]{mathalfa}
\usepackage{xfrac}
\usepackage[unicode,hidelinks,bookmarks=false]{hyperref}
\newcommand{\ie}{i.e.\ }
\newcommand{\cf}{cf.\ }
\newcommand{\resp}{respectively\ }
\newcommand{\Subsection}{\subsection}
\providecommand{\nobreakdash}{\nobreak\hbox{-}}
\setlength{\emergencystretch}{2em}
\multlinegap0pt
\usepackage{amssymb}
\newtheorem{proposition}{Proposition}
\def \H{{\mathcal H}}
\def \R{{\mathbb R}}
\def \S{{\mathcal S}}
\title{The heat semigroup associated with the Jacobi--Cherednik operator and its applications}
\author{Anirudha Poria, Ramakrishnan Radha}
\date{Source: arXiv:2409.05376v2 (CC BY 4.0)}
\begin{document}
\maketitle

\noindent\emph{Source and licence.} The heat semigroup associated with the Jacobi--Cherednik operator and its applications. Version arXiv:2409.05376v2 (CC BY 4.0), distributed by arXiv under the Creative Commons Attribution 4.0 International licence, http://creativecommons.org/licenses/by/4.0/, as recorded in the arXiv licence metadata for this exact version and shown on the abstract page https://arxiv.org/abs/2409.05376v2. Full licence notice: \texttt{licenses/arxiv-2409.05376-CC-BY-4.0.txt}.\par\smallskip

\noindent\textit{Editorial scope.} Counted unit: the article's proposition pro3 (source0.tex lines 319--341). The article's complete original proof of it is delivered with it (source0.tex lines 342--373, one \texttt{proof} environment of 32 lines). No mathematics is written, repaired or re-derived: the statement and the proof are the article's own lines, in the article's own notation, and no retained line is substituted or re-flowed.  Retained uncounted as the source-local context the proof needs: lines 220--224; lines 297--299; lines 301--303; lines 311--313; lines 315--317.  Nothing was omitted from any retained span, so the package carries no omission record.  No retained line contains a citation, so the wrapper carries no bibliography and no bibliography entry.  No retained line contains a figure environment, an image-inclusion command or drawing code, so no image is delivered and the package declares no asset dependency.  Editorial formatting only, in addition: an AMS \texttt{amsart} preamble that restates the article's own declarations and macros used by the retained lines, namely the environments \texttt{proposition} on separate counters, together with the standard packages those lines need.  The retained environments print in this document as Proposition 1 in order of appearance, whereas the article prints the same items with the numbering of its own full document.  The complete native source of the article as distributed in the arXiv e-print for this exact version is bundled unchanged as \texttt{sources/164970/source0.tex}.
\begin{equation}\label{eq02}
T^2_{\alpha, \beta} G^{\alpha, \beta}_\lambda(x)=i \lambda \; T_{\alpha, \beta} G^{\alpha, \beta}_\lambda(x)
=i \lambda \left( i \lambda G^{\alpha, \beta}_\lambda(x) \right)
=- \lambda^2  G^{\alpha, \beta}_\lambda(x).
\end{equation}

\begin{equation}\label{eq08}
\mathcal{F}_t^{\alpha, \beta}(x)= \H_{\alpha, \beta}^{-1}  \left(  e^{-\frac{t}{2}\left( \lambda^2 + \rho^2 \right)}  \right)(x), \quad x \in \R, \; t>0,
\end{equation}

\begin{proposition}\label{pro2}
For any $t>0$, the function $\mathcal{F}_t^{\alpha, \beta} \in \S^r_{\alpha, \beta} (\mathbb{R})$, for every $0< r \leq 1$. Further, the function $(x, t) \mapsto \mathcal{F}_t^{\alpha, \beta} (x)$, is infinitely differentiable on $ \mathbb{R} \times (0, \infty)$ and is a solution of $\mathrm{H} u = 0$.
\end{proposition} 

\begin{equation}\label{eq09}
\tau^{\alpha, \beta}_x f(y) = \int_{\mathbb{R}} \H_{\alpha, \beta} (f) (\lambda) \; G^{\alpha, \beta}_{-\lambda}(x) \; G^{\alpha, \beta}_{-\lambda}(y) \;  d\sigma_{\alpha, \beta}(\lambda), \quad \text{for all } y \in \R.
\end{equation}

\begin{equation}\label{eq1}
p_t^{\alpha, \beta}(x, y) =  \tau^{\alpha, \beta}_{-x} \mathcal{F}_t^{\alpha, \beta}(y).
\end{equation}

\begin{proposition}\label{pro3}
The Jacobi--Cherednik heat kernel $p_t^{\alpha, \beta}$ satisfies the following properties:
\begin{itemize}
\item[(a)] For all $x, y \in \R$ and $t>0$, we have
\begin{equation}\label{eq3}
p_t^{\alpha, \beta}(x, y)= \int_{\mathbb{R}} e^{-\frac{t}{2}\left( \lambda^2+\rho^2\right)}  G^{\alpha, \beta}_\lambda(x)\;  G^{\alpha, \beta}_\lambda(-y) \; d\sigma_{\alpha, \beta}(\lambda).
\end{equation}

\item[(b)] For all $x \in \R$ and $t>0$, $p_t^{\alpha, \beta}(x, 0) = \mathcal{F}_t^{\alpha, \beta}(x)$. 

\item[(c)] For all $x, y \in \R$ and $t>0$, $p_t^{\alpha, \beta}(x, y)=p_t^{\alpha, \beta}(-y, -x)$.

\item[(d)] For all $x \in \R$, $t>0$ and $\lambda \in \mathbb{C}$, we have $p_t^{\alpha, \beta}(x, \cdot) \in \S^r_{\alpha, \beta} (\mathbb{R})$ and 
\[ \H_{\alpha, \beta} \left( p_t^{\alpha, \beta}(x, - \; \cdot) \right) (\lambda) = e^{-\frac{t}{2}\left( \lambda^2+\rho^2\right)} G^{\alpha, \beta}_\lambda(x). \] 

\item[(e)] For all $x \in \R$ and $t>0$, we have
\begin{equation}\label{eq4}
\int_{\R} p_t^{\alpha, \beta}(x, y) A_{\alpha, \beta} (y) dy = e^{-\frac{t}{2} \rho^2}.
\end{equation}

\item[(f)] For fixed $y \in \R$, the function $u(x, t) = p_t^{\alpha, \beta}(x, y)$ solves the Jacobi--Cherednik heat equation $\mathrm{H} u=0$ on $\R \times (0, \infty)$.
\end{itemize}
\end{proposition}

\begin{proof}
(a) Since $G^{\alpha, \beta}_{\lambda}(\mu x)=G^{\alpha, \beta}_{\lambda \mu}(x)$, for $x, \lambda, \mu \in \R$, using (\ref{eq08}) and (\ref{eq09}), we get
\begin{eqnarray*}
p_t^{\alpha, \beta}(x, y) =  \tau^{\alpha, \beta}_{-x} \mathcal{F}_t^{\alpha, \beta}(y)
& = & \int_{\mathbb{R}} \H_{\alpha, \beta} \left(\mathcal{F}_t^{\alpha, \beta} \right) (\lambda) \; G^{\alpha, \beta}_{-\lambda}(-x) \; G^{\alpha, \beta}_{-\lambda}(y) \;  d\sigma_{\alpha, \beta}(\lambda) \\
& = & \int_{\mathbb{R}} e^{-\frac{t}{2}\left( \lambda^2+\rho^2\right)}  G^{\alpha, \beta}_{\lambda}(x) \; G^{\alpha, \beta}_{\lambda}(-y) \;  d\sigma_{\alpha, \beta}(\lambda).
\end{eqnarray*}
(b) Since $G^{\alpha, \beta}_\lambda(0)=1$, we obtain
\begin{eqnarray*}
p_t^{\alpha, \beta}(x, 0) =  \tau^{\alpha, \beta}_{-x} \mathcal{F}_t^{\alpha, \beta}(0)
& = & \int_{\mathbb{R}} e^{-\frac{t}{2}\left( \lambda^2+\rho^2\right)}  G^{\alpha, \beta}_{\lambda}(x) \; G^{\alpha, \beta}_{\lambda}(0) \;  d\sigma_{\alpha, \beta}(\lambda) \\
& = & \int_{\mathbb{R}} e^{-\frac{t}{2}\left( \lambda^2+\rho^2\right)}  G^{\alpha, \beta}_{\lambda}(x) \;  d\sigma_{\alpha, \beta}(\lambda) 
= \mathcal{F}_t^{\alpha, \beta}(x).
\end{eqnarray*}
(c) We have
\begin{eqnarray*}
p_t^{\alpha, \beta}(x, y) 
& = & \int_{\mathbb{R}} e^{-\frac{t}{2}\left( \lambda^2+\rho^2\right)}  G^{\alpha, \beta}_{\lambda}(x) \; G^{\alpha, \beta}_{\lambda}(-y) \;  d\sigma_{\alpha, \beta}(\lambda) \\
& = & \int_{\mathbb{R}} e^{-\frac{t}{2}\left( \lambda^2+\rho^2\right)}  G^{\alpha, \beta}_{\lambda}(-y) \;G^{\alpha, \beta}_{\lambda}(-(-x)) \;    d\sigma_{\alpha, \beta}(\lambda)
= p_t^{\alpha, \beta}(-y, -x).
\end{eqnarray*}
(d) From Proposition \ref{pro2}, we have $\mathcal{F}_t^{\alpha, \beta} \in \S^r_{\alpha, \beta} (\mathbb{R})$. Since the weighted Schwartz space $\S^r_{\alpha, \beta} (\mathbb{R})$ is invariant under these translation operators $\tau^{\alpha, \beta}_{x}$, using (\ref{eq1}), we obtain $p_t^{\alpha, \beta}(x, \cdot) \in \S^r_{\alpha, \beta} (\mathbb{R})$. Using the definition of $\H^{-1}_{\alpha, \beta}$, we get $\H^{-1}_{\alpha, \beta} \left( e^{-\frac{t}{2}\left( \lambda^2+\rho^2\right)} G^{\alpha, \beta}_\lambda(x) \right) (y)= p_t^{\alpha, \beta}(x, -y)$, and the second part follows immediately.\\
(e) From (d), we have 
\[ \H_{\alpha, \beta} \left( p_t^{\alpha, \beta}(x, - \; \cdot) \right) (\lambda) = e^{-\frac{t}{2}\left( \lambda^2+\rho^2\right)} G^{\alpha, \beta}_\lambda(x). \]
Therefore,
\[ \int_{\R} p_t^{\alpha, \beta}(x, -y) \; G^{\alpha, \beta}_\lambda(-y) \; A_{\alpha, \beta} (y) dy = e^{-\frac{t}{2}\left( \lambda^2+\rho^2\right)} G^{\alpha, \beta}_\lambda(x). \]
Using the change of variables, we get
\[\int_{\R} p_t^{\alpha, \beta}(x, y) \; G^{\alpha, \beta}_\lambda(y) \; A_{\alpha, \beta} (y) dy = e^{-\frac{t}{2}\left( \lambda^2+\rho^2\right)} G^{\alpha, \beta}_\lambda(x). \]
Since $G^{\alpha, \beta}_0(\cdot)=1$, by taking $\lambda=0$, we obtain
\[ \int_{\R} p_t^{\alpha, \beta}(x, y) A_{\alpha, \beta} (y) dy = e^{-\frac{t}{2} \rho^2}. \] 
(f) Finally, using differentiation under the integral sign in (a), interchanging the operator $T^2_{\alpha, \beta}$ with the integral sign, and using (\ref{eq02}), we get  $\mathrm{H} p_t^{\alpha, \beta}(x, y) = 0$.
\end{proof}

\end{document}
